% Recency Is Not Relevance: Certifying Provider-Native Context Manipulation
% arXiv-ready. Compiles on Overleaf (pdfLaTeX) with no external image files.
%
% SUBMISSION
%   Primary category : cs.SE  (Software Engineering)
%   Cross-list       : cs.LG  (Machine Learning)
%   Rationale: the contribution is a claim about how deployed agents behave when a
%   third party edits their context — a systems/engineering finding — and the method
%   is a test protocol, not a learning technique. cs.LG is the right cross-list
%   because the certification machinery (distribution-free bounds) reads as ML.
%
%   Upload BOTH files, or every reported number renders as "--":
%     docs/paper/provider_compaction.tex
%     docs/paper/generated/provider_compaction.tex   (the \input'd macro fragment)
%   Compiler: pdfLaTeX. See docs/PUBLISHING.md for the step-by-step.
%
% Every number in this paper is generated from the run artifacts by
%   python benchmarks/provider_compaction_to_latex.py
% which refuses to emit LaTeX unless each report's embedded protocol hash matches
% its protocol.json. Do not hand-edit numbers into the prose.

\documentclass[11pt]{article}

\usepackage[utf8]{inputenc}
\usepackage[T1]{fontenc}
\usepackage{lmodern}
\usepackage[margin=1in]{geometry}
\usepackage{amsmath}
\usepackage{amssymb}
\usepackage{booktabs}
\usepackage{graphicx}
\usepackage{xcolor}
\usepackage{url}
\usepackage[hidelinks]{hyperref}
\usepackage{cleveref}  % must load after hyperref

% Generated macros. If the fragment is missing, fail loudly rather than silently
% shipping a paper with placeholder numbers — the same guardrail main.tex uses.
\IfFileExists{generated/provider_compaction.tex}{%
  % Generated by benchmarks/provider_compaction_to_latex.py — do not edit.
% Source: benchmarks/results/provider-compaction/*/report.json

\newcommand{\pcAnthropicDefaultTwoAB}{100.0\%}
\newcommand{\pcAnthropicDefaultTwoABCI}{[+91.2, +100.0]}
\newcommand{\pcAnthropicDefaultTwoAA}{2.5\%}
\newcommand{\pcAnthropicDefaultTwoExcess}{+97.5pp}
\newcommand{\pcAnthropicDefaultTwoExcessCI}{[+92.5, +100.0]}
\newcommand{\pcAnthropicDefaultTwoBound}{100.0\%}
\newcommand{\pcAnthropicDefaultTwoFired}{40/40}
\newcommand{\pcAnthropicDefaultTwoCalls}{360}
\newcommand{\pcAnthropicDefaultTwoModel}{claude-opus-4-8}
\newcommand{\pcAnthropicDefaultOneAB}{95.0\%}
\newcommand{\pcAnthropicDefaultOneABCI}{[+83.5, +98.6]}
\newcommand{\pcAnthropicDefaultOneAA}{2.5\%}
\newcommand{\pcAnthropicDefaultOneExcess}{+92.5pp}
\newcommand{\pcAnthropicDefaultOneExcessCI}{[+85.0, +100.0]}
\newcommand{\pcAnthropicDefaultOneBound}{99.5\%}
\newcommand{\pcAnthropicDefaultOneFired}{40/40}
\newcommand{\pcAnthropicDefaultOneCalls}{360}
\newcommand{\pcAnthropicDefaultOneModel}{claude-opus-4-8}
\newcommand{\pcAnthropicAggressiveAB}{92.5\%}
\newcommand{\pcAnthropicAggressiveABCI}{[+80.1, +97.4]}
\newcommand{\pcAnthropicAggressiveAA}{7.5\%}
\newcommand{\pcAnthropicAggressiveExcess}{+85.0pp}
\newcommand{\pcAnthropicAggressiveExcessCI}{[+70.0, +95.0]}
\newcommand{\pcAnthropicAggressiveBound}{98.6\%}
\newcommand{\pcAnthropicAggressiveFired}{40/40}
\newcommand{\pcAnthropicAggressiveCalls}{360}
\newcommand{\pcAnthropicAggressiveModel}{claude-opus-4-8}
\newcommand{\pcOpenAIThreeKAB}{20.0\%}
\newcommand{\pcOpenAIThreeKABCI}{[+10.5, +34.8]}
\newcommand{\pcOpenAIThreeKAA}{2.5\%}
\newcommand{\pcOpenAIThreeKExcess}{+17.5pp}
\newcommand{\pcOpenAIThreeKExcessCI}{[+7.5, +30.0]}
\newcommand{\pcOpenAIThreeKBound}{36.7\%}
\newcommand{\pcOpenAIThreeKFired}{40/40}
\newcommand{\pcOpenAIThreeKCalls}{360}
\newcommand{\pcOpenAIThreeKModel}{gpt-5.2}
\newcommand{\pcOpenAIOneKAB}{12.5\%}
\newcommand{\pcOpenAIOneKABCI}{[+5.5, +26.1]}
\newcommand{\pcOpenAIOneKAA}{0.0\%}
\newcommand{\pcOpenAIOneKExcess}{+12.5pp}
\newcommand{\pcOpenAIOneKExcessCI}{[+2.5, +22.5]}
\newcommand{\pcOpenAIOneKBound}{27.8\%}
\newcommand{\pcOpenAIOneKFired}{40/40}
\newcommand{\pcOpenAIOneKCalls}{360}
\newcommand{\pcOpenAIOneKModel}{gpt-5.2}
\newcommand{\pcN}{40}
\newcommand{\pcVotes}{3}
\newcommand{\pcAlpha}{0.1}
\newcommand{\pcDelta}{0.05}
\newcommand{\pcAnthropicDefaultTwoCleared}{141,321}
\newcommand{\pcAnthropicDefaultOneCleared}{141,321}
\newcommand{\pcAnthropicAggressiveCleared}{94,211}

\newcommand{\pcTableBody}{%
Anthropic context editing --- \textbf{default}, run 2 & keep=3, trigger 3k & 100.0\% & 2.5\% & +97.5pp & 100.0\% & No \\
Anthropic context editing --- \textbf{default}, run 1 & keep=3, trigger 3k & 95.0\% & 2.5\% & +92.5pp & 99.5\% & No \\
Anthropic context editing --- aggressive & keep=0, trigger 500 & 92.5\% & 7.5\% & +85.0pp & 98.6\% & No \\
OpenAI server-side compaction & threshold 3k & 20.0\% & 2.5\% & +17.5pp & 36.7\% & No \\
OpenAI server-side compaction & threshold 1k & 12.5\% & 0.0\% & +12.5pp & 27.8\% & No \\
}
}{%
  \GenericWarning{}{WARNING: generated/provider_compaction.tex missing --- every
  reported number will render as a placeholder. Run
  benchmarks/provider_compaction_to_latex.py before distributing this PDF.}%
  \newcommand{\pcAnthropicDefaultTwoAB}{--}\newcommand{\pcAnthropicDefaultTwoAA}{--}%
  \newcommand{\pcAnthropicDefaultTwoAdj}{--}\newcommand{\pcAnthropicDefaultTwoBound}{--}%
  \newcommand{\pcAnthropicDefaultTwoFired}{--}\newcommand{\pcAnthropicDefaultTwoCalls}{--}%
  \newcommand{\pcAnthropicDefaultTwoModel}{--}%
  \newcommand{\pcAnthropicDefaultOneAB}{--}\newcommand{\pcAnthropicDefaultOneAA}{--}%
  \newcommand{\pcAnthropicDefaultOneAdj}{--}\newcommand{\pcAnthropicDefaultOneBound}{--}%
  \newcommand{\pcAnthropicDefaultOneFired}{--}\newcommand{\pcAnthropicDefaultOneCalls}{--}%
  \newcommand{\pcAnthropicDefaultOneModel}{--}%
  \newcommand{\pcAnthropicAggressiveAB}{--}\newcommand{\pcAnthropicAggressiveAA}{--}%
  \newcommand{\pcAnthropicAggressiveAdj}{--}\newcommand{\pcAnthropicAggressiveBound}{--}%
  \newcommand{\pcAnthropicAggressiveFired}{--}\newcommand{\pcAnthropicAggressiveCalls}{--}%
  \newcommand{\pcAnthropicAggressiveModel}{--}%
  \newcommand{\pcOpenAIThreeKAB}{--}\newcommand{\pcOpenAIThreeKAA}{--}%
  \newcommand{\pcOpenAIThreeKAdj}{--}\newcommand{\pcOpenAIThreeKBound}{--}%
  \newcommand{\pcOpenAIThreeKFired}{--}\newcommand{\pcOpenAIThreeKCalls}{--}%
  \newcommand{\pcOpenAIThreeKModel}{--}%
  \newcommand{\pcOpenAIOneKAB}{--}\newcommand{\pcOpenAIOneKAA}{--}%
  \newcommand{\pcOpenAIOneKAdj}{--}\newcommand{\pcOpenAIOneKBound}{--}%
  \newcommand{\pcOpenAIOneKFired}{--}\newcommand{\pcOpenAIOneKCalls}{--}%
  \newcommand{\pcOpenAIOneKModel}{--}%
  \newcommand{\pcN}{--}\newcommand{\pcVotes}{--}\newcommand{\pcAlpha}{--}%
  \newcommand{\pcDelta}{--}\newcommand{\pcAnthropicAggressiveCleared}{--}%
  \newcommand{\pcTableBody}{\multicolumn{7}{c}{\textit{macros missing}} \\}}

\title{\textbf{Recency Is Not Relevance:\\
Certifying Provider-Native Context Manipulation in LLM Agents}}

\author{%
  Chandra Shekhar Mudarapu\\
  \normalsize Independent Researcher\\
  \normalsize \texttt{chandu1221@gmail.com}%
  \thanks{Code, run artifacts, and the pre-registered protocols:
  \url{https://github.com/dshakes/distil}}%
}

\date{}

\begin{document}
\maketitle

\begin{abstract}
Both major LLM providers now rewrite the agent's context window on the server, by
default: Anthropic's \emph{context editing} deletes old tool results, and OpenAI's
\emph{server-side compaction} replaces them with a summary. Neither publishes what
this does to an agent's decisions, and neither has an incentive to --- no vendor
ships an instrument that measures its own degradation. We apply a
decision-equivalence test to both features under a pre-registered protocol with an
A/A self-agreement arm, majority-of-\pcVotes{} voting, and distribution-free upper
bounds. On a decision-bearing corpus of $n=\pcN$ tool-use episodes, Anthropic's
\emph{default} clearing policy (retain the 3 most recent tool uses) changed the
agent's next action in \pcAnthropicDefaultOneAB{} of cases, and
\pcAnthropicDefaultTwoAB{} on an independent replication of the same
pre-registered protocol, against an A/A noise floor of
\pcAnthropicDefaultOneAA{}. OpenAI's compaction changed \pcOpenAIThreeKAB{} under
matched conditions. Our central finding is mechanistic rather than numerical:
retaining the most recent $k$ tool results did \emph{not} reduce the
decision-change rate relative to deleting all of them
(\pcAnthropicAggressiveAB{} $\rightarrow$
\pcAnthropicDefaultOneAB{}--\pcAnthropicDefaultTwoAB{}); it changed the failure
\emph{mode}, converting visible stalls into confident wrong actions. Retention
counted in \emph{recency} cannot protect a decision that depends on
\emph{relevance}, and in agent trajectories the load-bearing facts are gathered
early and then buried under routine calls. We argue that the default policy is
therefore the more dangerous configuration, that summarization is empirically
$5$--$7\times$ safer than deletion, and that context manipulation should ship with
a decision-equivalence certificate rather than a token-reduction number.
\end{abstract}

\section{Introduction}
\label{sec:intro}

An LLM agent's context window is no longer under the application's sole control.
Anthropic's context editing~\cite{anthropiccontextediting} removes tool results
once the transcript crosses a token threshold; OpenAI's server-side
compaction~\cite{openaicompaction} summarizes prior turns into an opaque item. Both are enabled through a single request parameter, both operate
server-side, and both are promoted as the recommended path for long-running
agents. The user-visible contract is about \emph{tokens}: fewer of them, and
therefore lower cost and longer sessions.

The contract that matters, however, is about \emph{behavior}. An agent is a
sequence of decisions, each conditioned on what remains in the window. If removing
or summarizing a tool result changes which tool the agent calls next, the token
saving has been paid for in a currency the vendor does not report.

We make that currency measurable. Our contribution is not a compression method but
an instrument and the result of pointing it at the two features that now sit,
by default, between every agent and its model.

\paragraph{Contributions.}
\begin{enumerate}
  \item A pre-registered decision-equivalence protocol for third-party
    certification of context manipulation, with an A/A arm that separates true
    behavioral change from model sampling noise, ground-truthed firing detection,
    and a distribution-free upper bound on the change rate
    (\Cref{sec:method}).
  \item The first published measurement of what provider-native context
    manipulation does to agent decisions, for both major providers, at both
    aggressive and default configurations (\Cref{sec:results}).
  \item A mechanistic explanation of why the \emph{default} retention policy is
    more dangerous than deleting everything: recency-based retention cannot
    protect relevance-based dependencies, and it converts loud failures into
    silent ones (\Cref{sec:mechanism}).
\end{enumerate}

\section{Related work}
\label{sec:related}

\paragraph{Prompt and context compression.}
Work on prompt compression has largely optimized a reconstruction- or
task-accuracy objective: LLMLingua~\cite{llmlingua,longllmlingua,llmlinguatwo} and
its successors select tokens by
perplexity-based salience, while later methods compress with learned encoders or
structured intermediate representations. This literature evaluates compression by
what is preserved in the \emph{text}, or by end-task accuracy on static
benchmarks.

\paragraph{Compaction for long-horizon agents.}
Recent work moves the question to agent trajectories. ACON~\cite{acon} optimizes
natural-language compression guidelines by contrasting successful full-context
trajectories against failed compressed ones. Work on execution
instability~\cite{trace} argues
that repeatedly replacing working context with summaries causes agents to lose
track of execution state, and evaluates compression at the \emph{boundary} where
it occurs through paired continuations rather than at terminal outcomes. Studies of
implicit compression for software-engineering agents~\cite{implicitswe} report
that reconstruction errors --- hallucinated file paths, altered line numbers --- break tools that
require exact matching. A complementary line replaces token-level salience with
typed, decision-oriented fields, retaining what a unit implies for the next action
rather than which of its tokens score highly~\cite{damc} --- a design that is
explicitly relevance-indexed rather than recency-indexed, and therefore not subject
to the failure we characterize. Our finding is consistent with this line and extends it in
two ways: we measure a \emph{deployed, default-on} provider feature rather than a
research method, and we identify retention-by-recency as a specific mechanism by
which a policy that appears conservative fails.

\paragraph{Certification and risk control.}
Distribution-free risk control~\cite{crc,ltt} has been applied to LLM outputs,
routing, and tool calls~\cite{rolestrat}.

\paragraph{Compression as an attack surface.}
Independently of accidental information loss, an adversary who knows a compressor
sits in the path can craft input that survives compression while the guardrails
around it are budgeted away~\cite{coma}. That work targets compressors under the
application's control; the features we study are enabled provider-side, which
removes even the option of auditing the transformation locally. The two concerns
compound: a manipulation that changes decisions on benign input is also a
manipulation whose behavior on adversarial input is unmeasured. We use the same machinery for a different object: not the model's output,
but a transformation applied to the model's \emph{input} by a third party.

\section{Method}
\label{sec:method}

\paragraph{Decision-equivalence.}
Fix an episode: a multi-turn transcript ending at a point where the agent must act.
Let $A$ denote the transcript with manipulation disabled and $B$ the same
transcript with the provider feature enabled. The manipulation is
\emph{decision-preserving} on that episode iff the agent's next action is
unchanged. Actions are compared by a normalized signature over the tool name and
arguments, so that formatting differences do not register as behavioral change and
a switch from a tool call to prose does.

\paragraph{The A/A arm.}
A raw $A/B$ change rate is uninterpretable on its own, because a nonzero
temperature model does not always agree with itself. We therefore run a second,
independent baseline arm on identical input. The resulting $A/A$ rate is the
self-agreement floor; we report it alongside every $A/B$ rate and also report an
adjusted rate that discounts it. Each arm's decision is the majority of
\pcVotes{} samples.

\paragraph{Pre-registration.}
The protocol --- $n$, $\alpha$, $\delta$, vote count, trigger threshold, retention
parameter, and model --- is serialized to \texttt{protocol.json} before the first
API call, and its SHA-256 is embedded in the resulting report. Choosing $n$ after
seeing results is the failure mode this design exists to prevent. The LaTeX
generator for this paper recomputes that hash and refuses to emit numbers if the
pairing is broken.

\paragraph{Firing ground truth.}
An episode contributes to the sample only if the manipulation actually fired,
determined from the provider's own response metadata rather than inferred from
token counts. Episodes where the feature did not fire are excluded rather than
counted as successes, which would otherwise dilute the rate toward zero.

\paragraph{Bounds.}
We report a distribution-free upper confidence bound on the change rate at
$\delta=\pcDelta$, and declare a feature \emph{certified decision-safe} at
$\alpha=\pcAlpha$ only if that bound falls below $\alpha$.

\section{Experimental setup}
\label{sec:setup}

\paragraph{Corpora.}
Episodes are seeded, byte-reproducible synthetic tool-use transcripts in a
customer-support domain, in which every decision-bearing fact (order state,
billing record) appears inside a tool result surrounded by realistic log noise, and
no directive tells the model what to do. The first corpus uses two tool rounds; the
second uses seven, ordered so that the two decision-bearing rounds occur at the
head and are followed by four routine rounds (inventory, shipping scans,
communications log, promotions) carrying no decision-relevant fact.

\paragraph{Configurations.}
For Anthropic we test an aggressive setting (retain zero tool uses, low trigger)
and the shipped \emph{default} retention (retain the 3 most recent tool uses). For
OpenAI we test two compaction thresholds. Live API calls are hard-capped in code
per run; each run made \pcAnthropicDefaultOneCalls{} calls. Subjects are
\texttt{\pcAnthropicDefaultOneModel} and \texttt{\pcOpenAIOneKModel}.

\section{Results}
\label{sec:results}

\begin{table}[t]
\centering
\footnotesize
\setlength{\tabcolsep}{3pt}
\begin{tabular}{llccccc}
\toprule
Manipulation & Config & A/B & A/A & Adj. & 95\% bd. & Cert. \\
\midrule
\pcTableBody
\bottomrule
\end{tabular}
\caption{Decision-change rates for provider-native context manipulation,
$n=\pcN$ fired episodes per run, majority-of-\pcVotes{} decisions per arm.
``Certified'' is at $\alpha=\pcAlpha$ against the distribution-free upper bound.
No configuration of either feature certifies decision-safe.}
\label{tab:main}
\end{table}

\Cref{tab:main} reports every run. Three results stand out.

\paragraph{The default policy is not safer than deleting everything.}
Anthropic's aggressive configuration changed \pcAnthropicAggressiveAB{} of
decisions. Its \emph{default} retention changed \pcAnthropicDefaultOneAB{} and
\pcAnthropicDefaultTwoAB{} across two executions of the same pre-registered
protocol. Retaining the three most recent tool uses bought no reduction in
decision-change rate at all.

\paragraph{Summarization is safer than deletion.}
OpenAI's compaction changed \pcOpenAIOneKAB{} and \pcOpenAIThreeKAB{} of decisions
at the two thresholds tested, against \pcAnthropicAggressiveAB{} and
\pcAnthropicDefaultOneAB{}--\pcAnthropicDefaultTwoAB{} for deletion --- a gap of
roughly $5$--$7\times$ in the same harness, on the same corpora, with the same
grader.

\paragraph{Nothing certifies.}
No configuration of either feature has an upper bound below $\alpha=\pcAlpha$.
Both numbers are published regardless of which direction they point; selecting the
kinder run after the fact is the practice this protocol exists to refuse, which is
why both Anthropic default-policy executions appear in \Cref{tab:main}.

\section{Mechanism: recency is not relevance}
\label{sec:mechanism}

The rates alone understate the finding. The \emph{shape} of the failures differs
between configurations in a way that matters more than their magnitude.

Under the aggressive configuration, with all tool results cleared, the dominant
failure was the agent ceasing to act: the large majority of flips were
tool-call$\rightarrow$prose. Having lost the records it needed, the agent produced
text instead of an action. \pcAnthropicAggressiveCleared{} input tokens were cleared across the
episode set, as reported by the provider's own metadata.

Under the default retention policy, those stalls largely disappeared and were
replaced by tool-call$\rightarrow$\emph{different}-tool flips. The reason is
structural. The three retained results are the three \emph{most recent} ones ---
in our second corpus, the routine rounds. They are sufficient to keep the agent in
an acting posture: the window still looks populated, recent, and coherent. What is
gone is the older, decision-bearing evidence. The agent therefore acts confidently
on an incomplete record, issuing a refund where the policy requires a fraud
escalation.

This is the central claim of the paper, and it is not specific to one vendor's
implementation:

\begin{quote}
A retention rule indexed by \emph{recency} cannot protect a decision that depends
on \emph{relevance}. In agent trajectories, evidence is typically gathered early
and then buried under routine follow-up calls, so recency-ordered retention
preferentially discards exactly the load-bearing facts.
\end{quote}

The safety consequence is asymmetric. A stalled agent is a visible incident: a
human sees prose where an action belonged. An agent that acts on a partially
erased record produces a plausible, well-formed, wrong action --- in an agent with
write access, a silent incorrect write. The default policy converts the first
failure mode into the second while leaving the total failure rate unchanged. By
the metric a vendor is likely to track (does the agent still produce tool calls?)
the default looks like an improvement.

\section{Threats to validity}
\label{sec:threats}

\paragraph{Synthetic corpora, deliberately ordered.}
Both corpora are synthetic, and the second is constructed so that recency-ordered
retention keeps only routine rounds. This measures what the policy's mechanism does
when load-bearing facts age out of the retention window; it is \emph{not} a
prevalence estimate for how often that ordering arises in production traffic. We
claim the mechanism, not the base rate.

\paragraph{Lowered trigger.}
Experiments hold the real default retention rule but fire it at a reduced token
threshold, because episodes sized for a bounded API budget cannot cross the
production default. What generalizes is the retention rule; the threshold is a
lever to make it fire affordably. The behavior of the default threshold on genuine
long-horizon transcripts remains open and requires proportionally larger budget.

\paragraph{Decision-equivalence is not task success.}
The measured invariant is the agent's next action. Per-step equivalence and
end-to-end task success can diverge in both directions, and a change in next
action is not automatically a task failure.

\paragraph{Point-in-time.}
Both features are evolving surfaces measured on the dates and model versions
recorded in the artifacts. Vendors may change behavior without notice; these
results should be reproduced, not extrapolated.

\section{Conclusion}
\label{sec:conclusion}

Context manipulation has become infrastructure: enabled by default, executed
server-side, and reported in tokens. We show that the reporting unit is wrong.
Anthropic's default context-editing policy changed the agent's next action in
\pcAnthropicDefaultOneAB{}--\pcAnthropicDefaultTwoAB{} of pre-registered,
replicated cases, and its retention parameter --- the part that looks like a safety
margin --- bought no reduction in that rate while converting visible failures into
silent ones. Summarization fared substantially better than deletion, but neither
certifies decision-safe at $\alpha=\pcAlpha$.

The remedy is not to disable these features; it is to report them honestly. A
context operation should ship with a decision-equivalence certificate the way a
model ships with an evaluation. The harness used here is open source and runs
against either provider on a user's own configuration and traffic.

\paragraph{Reproducibility.}
The corpora are seeded and byte-reproducible; every run directory contains the
pre-registration written before the first API call and a report whose embedded
hash binds it to that pre-registration. The commands, the artifacts, and the
generator that turns them into every number in this paper are in the repository.

\begin{thebibliography}{99}

\bibitem{llmlingua} H.~Jiang, Q.~Wu, C.-Y. Lin, Y.~Yang, L.~Qiu.
\emph{LLMLingua: Compressing Prompts for Accelerated Inference of Large Language
Models.} EMNLP 2023. arXiv:2310.05736.

\bibitem{longllmlingua} H.~Jiang, Q.~Wu, X.~Luo, D.~Li, C.-Y. Lin, Y.~Yang, L.~Qiu.
\emph{LongLLMLingua: Accelerating and Enhancing LLMs in Long Context Scenarios via
Prompt Compression.} ACL 2024. arXiv:2310.06839.

\bibitem{llmlinguatwo} Z.~Pan, Q.~Wu, H.~Jiang, M.~Xia, X.~Luo, J.~Zhang, Q.~Lin,
V.~R\"uhle, Y.~Yang, C.-Y. Lin, H.~V. Zhao, L.~Qiu, D.~Zhang.
\emph{LLMLingua-2: Data Distillation for Efficient and Faithful Task-Agnostic
Prompt Compression.} Findings of ACL 2024. arXiv:2403.12968.

\bibitem{acon} M.~Kang, W.-N. Chen, D.~Han, H.~A. Inan, L.~Wutschitz, Y.~Chen,
R.~Sim, S.~Rajmohan.
\emph{ACON: Optimizing Context Compression for Long-horizon LLM Agents.}
ICML 2026. arXiv:2510.00615.

\bibitem{trace} G.~Min, L.~Wu, M.~Darbari, C.~Chen, L.~Hong.
\emph{Toward Reliable Context Compression for Long-Horizon Agents: An Empirical
Study of Execution Instability.} arXiv:2608.06503, 2026.

\bibitem{damc} X.~Guan, Q.~Zhao, Y.~Deng.
\emph{Decision-Aware Memory Cards: Counterfactual-Inspired Context Selection and
Compression for Tool-Using LLM Agents.} arXiv:2606.08151, 2026.

\bibitem{implicitswe} K.~Gelvan, I.~Slinko, F.~Steinbauer, E.~Bogomolov,
F.~Kofler, Y.~Zharov.
\emph{On Problems of Implicit Context Compression for Software Engineering
Agents.} arXiv:2605.11051, 2026.

\bibitem{coma} Z.~Liu, Z.~Zhang, Y.~Xie, D.~She.
\emph{When Compression Becomes an Attack Surface: Black-Box Attacks on
Prompt-Compressed LLM Agents.} arXiv:2510.22963, 2026.

\bibitem{crc} A.~N. Angelopoulos, S.~Bates, A.~Fisch, L.~Lei, T.~Schuster.
\emph{Conformal Risk Control.} ICLR 2024. arXiv:2208.02814.

\bibitem{ltt} A.~N. Angelopoulos, S.~Bates, E.~Cand\`es, M.~Jordan, L.~Lei.
\emph{Learn Then Test: Calibrating Predictive Algorithms to Achieve Risk Control.}
Annals of Applied Statistics, 2025. arXiv:2110.01052.

\bibitem{rolestrat} M.~A. Rahman, M.~A. Rahman, N.~H. Samin, K.~R. Tasnia,
M.~H. Amin, S.~R. Ahona, J.~A. Noshin.
\emph{Beyond Aggregate Risk: Role-Stratified Conformal Risk Control for LLM Tool
Calls.} arXiv:2607.24343, 2026.

\bibitem{anthropiccontextediting} Anthropic.
\emph{Context editing.} Claude Platform Documentation.
\url{https://platform.claude.com/docs/en/build-with-claude/context-editing}.
Accessed 2026-08.

\bibitem{openaicompaction} OpenAI.
\emph{Compaction.} OpenAI API Documentation.
\url{https://developers.openai.com/api/docs/guides/compaction}.
Accessed 2026-08.

\end{thebibliography}

\end{document}
